% =========================================================
% Section: Mean Value Theorem
% =========================================================
% ---------------------------------------------------------

\begin{remark*}[Interpretation]
Rolle's theorem says that a differentiable curve returning to the same height
must level off somewhere between. The Mean Value Theorem tilts that picture:
between any two points, some interior tangent runs parallel to the chord.
Cauchy's version compares two functions at once and later powers
L'H\^{o}pital's rule.
\end{remark*}

\begin{figure}[H]
\centering
\resizebox{0.72\linewidth}{!}{%
\begin{tikzpicture}[scale=0.82]
  \def\h{1.5*sin(72*\x)}
  \glowcurve{color=atlasred, domain=0:5, axis=-1.7, top=1.9, expr={\h}}
  \draw[atlas axes] (0,0)--(5.4,0);
  \draw[atlas axes] (0,-1.9)--(0,2.1);
  \pgfmathsetmacro\xmax{1.25}
  \pgfmathsetmacro\xmin{3.75}
  \draw[atlas probe,atlasred] (\xmax-0.9,1.5)--(\xmax+0.9,1.5);
  \draw[atlas probe,atlasred] (\xmin-0.9,-1.5)--(\xmin+0.9,-1.5);
  \node[atlas dot] at (\xmax,1.5){};
  \node[atlas dot] at (\xmin,-1.5){};
  \node[atlas dot] at (0,0){};
  \node[atlas dot] at (5,0){};
  \node[atlas label,above] at (\xmax,1.62){\footnotesize horizontal tangent};
  \node[atlas label,below] at (\xmin,-1.62){\footnotesize horizontal tangent};
  \node[atlas label] at (3.95,0.62){\footnotesize$y=f(x)$};
  \node[atlas label,below left=-1pt] at (0,0){\footnotesize$a$};
  \node[atlas label,below right=-1pt] at (5,0){\footnotesize$b$};
  \node[atlas label,below] at (2.5,-1.95){\scriptsize$a<c<b$};
  \node[atlas label] at (2.25,1.2){\scriptsize$f(a)=f(b)$};
  \node[atlas label] at (2.55,0.55){\scriptsize$f'(c)=0$};
\end{tikzpicture}%
}
\caption{Rolle's Theorem: equal endpoint heights force a horizontal tangent.}
\label{fig:rolle}
\end{figure}


\begin{theorembox}{Theorem (Rolle's Theorem)}
\begin{theorem}[Rolle's Theorem]
\label{thm:rolles-theorem}
Let $a, b \in \mathbb{R}$ with $a < b$. Suppose that $f : [a,b] \to \mathbb{R}$ is continuous on $[a,b]$, differentiable on $(a,b)$, and satisfies $f(a) = f(b)$. Then there exists $c \in (a,b)$ such that
\[
f'(c) = 0.
\]
\hyperref[prf:rolles-theorem]{Go to proof.}
\end{theorem}
\end{theorembox}

\begin{remark*}[Standard quantified statement]
\begin{align*}
&\operatorname{IsContinuous}(f,[a,b])
\;\land\;
\forall x \in (a,b)\;\operatorname{IsDifferentiable}(f,x)
\;\land\;
f(a)=f(b)\\
&\quad\Rightarrow\;
\exists c \in (a,b)\;\operatorname{Derivative}(f,c,0).
\end{align*}
\end{remark*}

\begin{remark*}[Predicate reading]
\begin{align*}
&\operatorname{IsContinuous}(f,[a,b])
\;\land\;
\operatorname{IsDifferentiable}(f,(a,b))
\;\land\;
f(a)=f(b)\\
&\quad\Longrightarrow\;
\exists c\in(a,b)\;:\;\operatorname{Derivative}(f,c,0).
\end{align*}
\end{remark*}

\begin{remark*}[Failure modes]
\begin{description}
  \item[Exposition.]
Each hypothesis is essential:
\begin{itemize}
  \item \textbf{Continuity fails:} Let $f(x) = \operatorname{sgn}(x - 1/2)$ on $[0,1]$. Then $f(0) = f(1) = -1$, but $f$ has a jump discontinuity at $x = 1/2$. The derivative does not exist anywhere, so there is no horizontal tangent.
  \item \textbf{Differentiability fails:} Let $f(x) = |x - 1/2|$ on $[0,1]$. Then $f(0) = f(1) = 1/2$, and $f$ is continuous, but $f$ has a corner at $x = 1/2$ where the derivative is undefined. For all $c \ne 1/2$, we have $f'(c) = \pm 1 \ne 0$.
  \item \textbf{Equal endpoints fail:} Let $f(x) = x$ on $[0,1]$. Then $f$ is continuous and differentiable with $f'(x) = 1 > 0$ everywhere, but $f(0) = 0 \ne 1 = f(1)$.
\end{itemize}
\end{description}
\end{remark*}

\begin{remark*}[Interpretation]
\[
\bigl(\forall c \in (a,b)\;:\; f'(c) \ne 0\bigr)
\;\Rightarrow\;
f \text{ not continuous on } [a,b]
\;\lor\;
f \text{ not differentiable on } (a,b)
\;\lor\;
f(a) \ne f(b).
\]
\end{remark*}

\begin{remark*}[Interpretation]
\begin{align*}
&\bigl(\forall c\in(a,b)\;:\;\neg\operatorname{Derivative}(f,c,0)\bigr)\\
&\quad\Longrightarrow\;
\neg\operatorname{IsContinuous}(f,[a,b])
\;\lor\;
\neg\operatorname{IsDifferentiable}(f,(a,b))
\;\lor\;
f(a)\ne f(b).
\end{align*}
\end{remark*}

\begin{remark*}[Interpretation]
If a differentiable function returns to the same value at both endpoints of an interval, the derivative must vanish somewhere in between. Geometrically: if the curve starts and ends at the same height, at some interior point the tangent line is horizontal. The proof combines two ingredients --- the Extreme Value Theorem guarantees that a continuous function on a closed bounded interval attains its maximum and minimum, and the Interior Extremum Theorem guarantees that at an interior extremum the derivative is zero. If both extrema occur at the endpoints then $f$ is constant and $f' \equiv 0$ throughout. Rolle's Theorem is the special case of the MVT where the secant slope is zero.
\end{remark*}

\begin{dependencies}
\begin{itemize}
  \item \hyperref[thm:interior-extremum-theorem]{Theorem: Interior Extremum Theorem}
  \item \hyperref[thm:extreme-value-theorem]{Theorem: Extreme Value Theorem}
  \item \hyperref[def:derivative-at-a-point]{Definition: Derivative at a Point}
\end{itemize}
\end{dependencies}

% ---------------------------------------------------------

\begin{figure}[H]
\centering
\resizebox{0.72\linewidth}{!}{%
\begin{tikzpicture}[scale=0.82]
  \def\f{0.16*\x*\x+0.20*\x+0.55}
  \glowcurve{color=atlasblue, domain=0.5:4.5, axis=0, top=4.6, expr={\f}}
  \draw[atlas axes] (0,0)--(5.2,0);
  \draw[atlas axes] (0,0)--(0,4.6);
  \pgfmathsetmacro\fa{0.16*0.5*0.5+0.20*0.5+0.55}
  \pgfmathsetmacro\fb{0.16*4.5*4.5+0.20*4.5+0.55}
  \pgfmathsetmacro\msec{(\fb-\fa)/(4.5-0.5)}
  \pgfmathsetmacro\cc{(\msec-0.20)/0.32}
  \pgfmathsetmacro\fc{0.16*\cc*\cc+0.20*\cc+0.55}
  \draw[atlas probe,atlasgreen] (0.5,\fa)--(4.5,\fb);
  \draw[atlas probe,atlasred] ({\cc-1.7},{\fc-\msec*1.7})--({\cc+1.7},{\fc+\msec*1.7});
  \draw[dropline] (0.5,\fa)--(0.5,0);
  \draw[dropline] (\cc,\fc)--(\cc,0);
  \draw[dropline] (4.5,\fb)--(4.5,0);
  \node[atlas dot] at (0.5,\fa){};
  \node[atlas dot] at (4.5,\fb){};
  \node[atlas dot] at (\cc,\fc){};
  \node[atlas label,left=1pt] at (0.5,\fa){\footnotesize$(a,f(a))$};
  \node[atlas label,above right=-1pt] at (4.5,\fb){\footnotesize$(b,f(b))$};
  \node[atlas label,above right] at (3.2,2.35){\footnotesize$y=f(x)$};
  \node[atlas label,below right=1pt] at (\cc,\fc){\scriptsize$f'(c)=\dfrac{f(b)-f(a)}{b-a}$};
  \node[atlas label,below] at (0.5,-0.05){\footnotesize$a$};
  \node[atlas label,below] at (2.5,-0.05){\scriptsize$a<c<b$};
  \node[atlas label,below] at (4.5,-0.05){\footnotesize$b$};
\end{tikzpicture}%
}
\caption{Mean Value Theorem: an interior tangent matches the secant slope.}
\label{fig:mvt}
\end{figure}


\begin{theorembox}{Theorem (Mean Value Theorem)}
\begin{theorem}[Mean Value Theorem]
\label{thm:mean-value-theorem}
Let $a, b \in \mathbb{R}$ with $a < b$. Suppose that $f : [a,b] \to \mathbb{R}$ is continuous on $[a,b]$ and differentiable on $(a,b)$. Then there exists $c \in (a,b)$ such that
\[
f'(c) = \frac{f(b) - f(a)}{b - a}.
\]
\hyperref[prf:mean-value-theorem]{Go to proof.}
\end{theorem}
\end{theorembox}

\begin{remark*}[Standard quantified statement]
\begin{align*}
&\operatorname{IsContinuous}(f,[a,b])
\;\land\;
\forall x \in (a,b)\;\operatorname{IsDifferentiable}(f,x)\\
&\quad\Rightarrow\;
\exists c \in (a,b)\;
\operatorname{Derivative}\!\left(f,\,c,\,\tfrac{f(b)-f(a)}{b-a}\right).
\end{align*}
\end{remark*}

\begin{remark*}[Predicate reading]
\begin{align*}
&\operatorname{IsContinuous}(f,[a,b])
\;\land\;
\operatorname{IsDifferentiable}(f,(a,b))\\
&\quad\Longrightarrow\;
\exists c\in(a,b)\;:\;
\operatorname{Derivative}\!\left(f,\,c,\,\tfrac{f(b)-f(a)}{b-a}\right).
\end{align*}
\end{remark*}

\begin{remark*}[Failure modes]
\begin{description}
  \item[Exposition.]
Each hypothesis is essential:
\begin{itemize}
  \item \textbf{Continuity fails:} Let $f(x) = \begin{cases} 0 & \text{if } x \in [0,1/2) \\ 1 & \text{if } x \in [1/2,1] \end{cases}$. Then $f$ has a jump discontinuity at $x = 1/2$, and the secant slope is $\tfrac{f(1)-f(0)}{1-0} = 1$. But $f'(x) = 0$ wherever the derivative exists, so no interior point achieves the average rate of change.
  \item \textbf{Differentiability fails:} Let $f(x) = |x|$ on $[-1,1]$. Then $f$ is continuous and the secant slope is $\tfrac{f(1)-f(-1)}{1-(-1)} = \tfrac{1-1}{2} = 0$. But $f'(x) = \pm 1$ everywhere $f$ is differentiable (i.e., for $x \ne 0$), so no interior point has derivative equal to the secant slope.
\end{itemize}
The MVT is an existence result. It guarantees some $c$ exists but gives no formula for it.
\end{description}
\end{remark*}

\begin{remark*}[Interpretation]
\[
\left(\forall c \in (a,b)\;:\; f'(c) \ne \frac{f(b)-f(a)}{b-a}\right)
\;\Rightarrow\;
f \text{ not continuous on } [a,b]
\;\lor\;
f \text{ not differentiable on } (a,b).
\]
\end{remark*}

\begin{remark*}[Interpretation]
\begin{align*}
&\left(\forall c\in(a,b)\;:\;\neg\operatorname{Derivative}\!\left(f,\,c,\,\tfrac{f(b)-f(a)}{b-a}\right)\right)\\
&\quad\Longrightarrow\;
\neg\operatorname{IsContinuous}(f,[a,b])
\;\lor\;
\neg\operatorname{IsDifferentiable}(f,(a,b)).
\end{align*}
\end{remark*}

\begin{remark*}[Interpretation]
The MVT asserts that the instantaneous rate of change equals the average rate of change at some interior point. Geometrically: the tangent line at some point $c$ is parallel to the secant line through $(a,f(a))$ and $(b,f(b))$. The proof reduces to Rolle's Theorem by subtracting the secant: define $g(x) = f(x) - \tfrac{f(b)-f(a)}{b-a}(x-a)$, which satisfies $g(a)=g(b)=f(a)$, then apply Rolle's. The theorem is an existence result, not a construction: it guarantees some $c$ exists but gives no formula for it. Its immediate corollaries --- the constant function corollary, the monotonicity corollary, and the Lipschitz bound corollary --- each flow from the MVT by a one-line argument and are the theorem's primary applications.
\end{remark*}

\begin{dependencies}
\begin{itemize}
  \item \hyperref[thm:rolles-theorem]{Theorem: Rolle's Theorem}
  \item \hyperref[def:derivative-at-a-point]{Definition: Derivative at a Point}
\end{itemize}
\end{dependencies}

% ---------------------------------------------------------

\begin{theorembox}{Theorem (Cauchy Mean Value Theorem)}
\begin{theorem}[Cauchy Mean Value Theorem]
\label{thm:cauchy-mean-value-theorem}
Let $a,b\in\mathbb{R}$ with $a<b$. Suppose that $f,g:[a,b]\to\mathbb{R}$ are continuous on $[a,b]$ and differentiable on $(a,b)$. Then there exists $c\in(a,b)$ such that
\[
  \bigl(f(b)-f(a)\bigr)g'(c)
  =
  \bigl(g(b)-g(a)\bigr)f'(c).
\]
If, in addition, $g'(x)\neq 0$ for every $x\in(a,b)$, then $g(b)\neq g(a)$ and
\[
  \frac{f'(c)}{g'(c)}
  =
  \frac{f(b)-f(a)}{g(b)-g(a)}.
\]
\hyperref[prf:cauchy-mean-value-theorem]{Go to proof.}
\end{theorem}
\end{theorembox}

\begin{remark*}[Standard quantified statement]
\begin{align*}
&\operatorname{IsContinuous}(f,[a,b])
\land \operatorname{IsContinuous}(g,[a,b])
\land \operatorname{IsDifferentiable}(f,(a,b))
\land \operatorname{IsDifferentiable}(g,(a,b))\\
&\quad\Longrightarrow\;
\exists c\in(a,b)\;
\bigl[(f(b)-f(a))g'(c)=(g(b)-g(a))f'(c)\bigr].
\end{align*}
\end{remark*}

\begin{remark*}[Predicate reading]
The Cauchy mean value statement asserts
\[
\exists c\in(a,b)\;
\bigl[(f(b)-f(a))g'(c)=(g(b)-g(a))f'(c)\bigr],
\]
the existence of an interior point where the velocity vector
\((g'(c),f'(c))\) is parallel to the chord vector
\((g(b)-g(a),f(b)-f(a))\). Taking \(g(x)=x\) recovers the ordinary Mean Value Theorem.
\end{remark*}

\begin{remark*}[Negated quantified statement]
\[
  \forall c\in(a,b):\ 
  \bigl(f(b)-f(a)\bigr)g'(c)\neq
  \bigl(g(b)-g(a)\bigr)f'(c).
\]
\end{remark*}

\begin{remark*}[Failure modes]
\begin{description}
  \item[Exposition.]
The continuity and differentiability hypotheses fail for the same reasons they fail in the ordinary MVT: jumps break endpoint-to-interior control, and corners leave no derivative to match. The quotient form has the additional nonvanishing hypothesis \(g'\neq 0\) on \((a,b)\); without it the cross-multiplied conclusion may still hold, but the displayed quotient can be undefined.
\end{description}
\end{remark*}

\begin{remark*}[Interpretation]
\[
\begin{aligned}
&\bigl(\forall c\in(a,b):\ (f(b)-f(a))g'(c)\neq(g(b)-g(a))f'(c)\bigr)\\
&\quad\Longrightarrow\
\neg\operatorname{IsContinuous}(f,[a,b])
\vee \neg\operatorname{IsContinuous}(g,[a,b])
\vee \neg\operatorname{IsDifferentiable}(f,(a,b))
\vee \neg\operatorname{IsDifferentiable}(g,(a,b)).
\end{aligned}
\]
\end{remark*}

\begin{remark*}[Interpretation]
If no interior point makes the derivative pair \((g'(c),f'(c))\) parallel to
the chord pair \((g(b)-g(a),f(b)-f(a))\), then at least one of the four
regularity hypotheses fails: continuity of \(f\) or \(g\) on \([a,b]\), or
differentiability of \(f\) or \(g\) on \((a,b)\).
\end{remark*}

\begin{remark*}[Interpretation]
This is the parametric Mean Value Theorem. For the planar curve \(t\mapsto(g(t),f(t))\), some interior tangent is parallel to the chord joining its endpoints. It compares two rates of change rather than comparing one function to the horizontal axis. This is the form that later underlies L'Hopital-type ratio arguments and the Cauchy form of Taylor's remainder.
\end{remark*}

\begin{dependencies}
\begin{itemize}
  \item \hyperref[thm:rolles-theorem]{Theorem: Rolle's Theorem}
  \item \hyperref[def:derivative-at-a-point]{Definition: Derivative at a Point}
  \item \hyperref[thm:sum-rule]{Theorem: Sum Rule}
  \item \hyperref[thm:constant-multiple-rule]{Theorem: Constant Multiple Rule}
\end{itemize}
\end{dependencies}

% ---------------------------------------------------------

\begin{corollarybox}{Corollary (Derivative Bound Implies Lipschitz)}
\begin{corollary}[Derivative Bound Implies Lipschitz]
\label{cor:derivative-bound-implies-lipschitz}
Let \(I\subseteq\mathbb{R}\) be an interval and let \(f:I\to\mathbb{R}\) be differentiable on \(I\). If there exists \(M\geq 0\) such that
\[
  |f'(x)|\leq M
  \qquad\text{for every }x\in I,
\]
then \(f\) is Lipschitz on \(I\) with constant \(M\):
\[
  |f(x)-f(y)|\leq M|x-y|
  \qquad\text{for every }x,y\in I.
\]
\hyperref[prf:derivative-bound-implies-lipschitz]{Go to proof.}
\end{corollary}
\end{corollarybox}

\begin{remark*}[Standard quantified statement]
\[
  \bigl(\exists M\geq 0\;\forall x\in I:\ |f'(x)|\leq M\bigr)
  \Longrightarrow
  \forall x,y\in I:\ |f(x)-f(y)|\leq M|x-y|.
\]
\end{remark*}

\begin{remark*}[Predicate reading]
\[
  \operatorname{IsBoundedDerivative}(f,I,M)
  \Longrightarrow
  \operatorname{IsLipschitz}(f,I,M).
\]
\end{remark*}

\begin{remark*}[Failure modes]
\begin{description}
  \item[Exposition.]
The derivative bound must hold throughout the interval. A function such as \(f(x)=\sqrt{x}\) on \([0,1]\) is continuous and uniformly continuous, but its derivative is unbounded near \(0\), so no global Lipschitz constant follows from this test. The interval hypothesis is also essential: across a disconnected gap, the MVT cannot compare the two components.
\end{description}
\end{remark*}

\begin{remark*}[Interpretation]
This corollary converts a pointwise slope bound into a global modulus of continuity. It is the precise bridge from derivative estimates to the Lipschitz condition: controlling the size of \(f'\) controls every secant slope on the interval.
\end{remark*}

\begin{dependencies}
\begin{itemize}
  \item \hyperref[thm:mean-value-theorem]{Theorem: Mean Value Theorem}
  \item \hyperref[def:lipschitz-condition]{Definition: Lipschitz Condition}
  \item \hyperref[def:derivative-at-a-point]{Definition: Derivative at a Point}
\end{itemize}
\end{dependencies}
